\chapter{第五讲\quad 客观要件一：行为}
\label{lecture:05}
\noindent{\small\hyperref[ref:lecture:05]{查阅本讲扩展规则}}\par

本讲先问行为改变了什么风险，再问谁依法承担阻断风险的义务。把人记为 $x$，法益记为 $g$，行为记为 $a$，应当采取的措施记为 $m$，待归属的结果记为 $r$。同一人可能在不同时间实施不同的行为，时间下标不可省略。

\section{第一节\quad 危害行为}

\paragraph{定义与教材模型}
用 $\mathrm{Ext}(a)$ 表示外部活动或有意义的静止。用 $\mathrm{Ctrl}(a)$ 表示行为受意识支配。用 $\mathrm{New}(a,g)$ 表示对法益 $g$ 创设或增加了待审查的危险。教材以有体性、有意性、有害性识别危害行为。本节的初筛记作
\[
\mathrm{Act}(a)=\mathrm{Ext}(a)\land\mathrm{Ctrl}(a),
\qquad
\mathrm{Haz}(a,g)=\mathrm{Act}(a)\land\mathrm{New}(a,g).
\]
$\mathrm{Ctrl}$ 指控制行为的能力，故意针对的是构成要件事实和结果。两者分层审查。刑法第十三条提供犯罪的实质判断背景，上式是教材模型，分则构成、风险许可和违法阻却仍须另行核查。

\paragraph{推导}
甲把即将被落石砸中头部的乙推开，石头只砸中乙的肩部。题设确认甲的动作只减弱同一落石机制的危害，未制造其他危险。事实给出 $\mathrm{Act}(a)$ 及 $\neg\mathrm{New}(a,g)$，所以
\[
\mathrm{Haz}(a,g)
   =1\land 0=0 .
\]
肩伤实际发生，仍不能从 $\mathrm{Injury}(r)$ 倒推 $\mathrm{New}(a,g)$。若甲为使婴儿逃离火灾，把婴儿从楼上抛下，坠落是新机制，即使总体损害预期下降，仍先认定 $\mathrm{New}$。然后依事实审查刑法第二十一条的紧急避险。可把两题的差别记为同一风险减弱和以新风险替代。

\paragraph{边界}
脑中希望他人死亡不等于外部行为。纯反射性举动可能有外观而缺乏行为控制。评价可回看行为人此前能够控制的独立行为，避免把先前过失与当时反射混成同一时点。另如仅劝成年人参加普通运动，题设无特殊风险，虽然劝者怀有恶意且对方受伤，仍缺 $\mathrm{New}$。

\paragraph{被害人自陷风险}
在教材的典型分流中，令 $S_y$ 表示被害人 $y$ 是风险的实际实行、支配者，$K_y$ 表示其能够认识该具体风险，$F_y$ 表示其能自由选择、控制或退出。则
\[
S_y\land K_y\land F_y
   \ \Longrightarrow\ \mathrm{SelfRisk}(y,r)
   \quad\text{（在该风险及其结果的范围内）}.
\]
若熟悉河况的成年人自行决定入水，且未受欺骗、强迫，也无他人特别保护义务的线索，在题设的通常自陷风险模型中，三个条件可以逐项认定。若改为甲自己掌舵，乙只同意冒险乘船，风险控制仍在甲，乙的 $S_y$ 不成立。此时继续审查甲制造的风险和主观罪过。同意冒险不等于有效承诺接受伤亡结果。紧急救助、信息不对称、依赖和实质支配都会改变结论。救援行为也可能正是先前侵害所引起的反应，应结合第七讲分析。

\section{第二节\quad 不作为犯的分类}

\paragraph{定义与规则层级}
把待审行为放回具体规范。$\mathrm{Ban}_j(a)$ 指罪名 $j$ 禁止实施的行为类型，$\mathrm{Cmd}_j(m)$ 指该规范命令实施的行为。以违反禁止规范的方式实现罪状是作为，以违反命令规范的方式实现罪状是不作为。这是教材分类。真正不作为犯的作为命令来自罪状本身，例如刑法第二百六十一条的遗弃罪。不真正不作为犯则以一般结果罪状评价特别义务人不作为，例如以不救助造成死亡来审查第二百三十二条的故意杀人罪。

\paragraph{推导}
设 $D$ 为特定扶养义务，$N$ 为不能独立生活，$O$ 为拒绝扶养，$B$ 为情节恶劣。取教材对第二百六十一条的四个核心条件，令 $\Gamma$ 记载负有扶养义务的甲有能力却持续拒绝给不能独立生活的乙提供生活保障，且已经达到情节恶劣。则
\[
\begin{aligned}
\Gamma&\vdash D\land N\land O\land B,\\
D\land N\land O\land B&\Longrightarrow
     \mathrm{Elements}_{261},\\
\Gamma&\vdash \mathrm{Elements}_{261}.
\end{aligned}
\]
$\mathrm{Elements}_{261}$ 表示该例待审的构成要件已经通过，归责还需主体、罪过与排除事由等前提。若甲以把乙移出住所的积极动作来拒绝扶养，外观虽有动作，评价的核心仍是应扶养而拒绝。改变动作外观不改变命令规范。

\paragraph{边界}
主动撤除救治设备并直接制造致死风险，可能同时有积极致害和不继续救助的两种描述。教材对同一过程可优先审查积极作为。只离开岗位而未救治则要先找特别义务。持有型犯罪通常按作为理解，以明知并实际控制且继续维持的状态为重点，例如刑法第三百四十八条还须审查毒品种类、数量等法定条件。一时接触未知物，随后发现并立即结束控制，与主动维持非法控制的事实结构不同。作为和不作为是客观行为形态，故意和过失是罪过形态，两组分类不能互相推出。

\section{第三节\quad 不作为犯的成立条件}

\paragraph{定义与教材规则}
对不真正不作为的既遂结果犯，设 $D(x,m,t)$ 为 $x$ 在时点 $t$ 对具体措施 $m$ 负有刑法意义上的特别义务，$C(x,m,t)$ 为当时可实施，$N(x,m,t)$ 为未实施，$A(m,r,t)$ 为依法可认定实施该措施具有阻止该结果的实效，$E_j$ 为该不作为与罪名 $j$ 的积极实行行为等价，$R(r)$ 为实际发生法定结果。先证明
\[
\mathrm{OmitObj}_j
 = E_j\land R(r)\land
   \exists m\,\bigl(D(x,m,t)\land C(x,m,t)\land
                  N(x,m,t)\land A(m,r,t)\bigr).
\]
存在量词要求同一个 $m$ 同时符合各条件。不能用「会报警」证明「能下水」，再用「下水能救」给报警补上因果性。$A$ 是对具体措施和结果的反事实证成，不能把任何微小的获救概率机械当作完成犯的证明。再接上 $j$ 所要求的故意或法定过失，并审查排除事由。真正不作为犯是否必须有结果，回到该罪自己的条文。

\paragraph{特别义务的两条路径}
教材的实质二分法记为 $D_{\text{源}}\lor D_{\text{人}}$。前者问谁监管危险物、负有监管关系的他人或者自己先前制造的危险。后者问法益为什么依赖此人保护。常见依据包括法律关系、职务业务、约定、自愿接管后形成的依赖，以及特定场所的控制与依赖。一般道德期待和一般报警义务，不能直接推出对一切损害结果的保证人义务。先行行为是否须违法、防卫人是否有救助义务，涉及学说分歧。应就被防卫者或第三人、允许的防卫结果与其后新生危险分别说明所采规则。

\paragraph{推导}
幼儿落入蓄水池，值守监护人甲知道幼儿处境，触手可及的救生杆足以救出幼儿，甲为使其死亡故意不伸杆，幼儿死亡。取 $m$ 为用杆拉救，则
\[
\begin{aligned}
\text{监护与依赖}&\Longrightarrow D(x,m,t),\\
\text{近处可用救生杆}&\Longrightarrow C(x,m,t),\\
\text{鉴定及现场条件}&\Longrightarrow A(m,r,t),\\
D\land C\land A\land N\land E_{232}\land R
   &\Longrightarrow\mathrm{OmitObj}_{232}.
\end{aligned}
\]
已证死亡故意，且题设排除阻却事由时，依第二百三十二条的教材模型得出不作为故意杀人既遂。若第三人及时救出，$R(r)$ 不成立。可再用第十一讲的着手和意志以外原因规则判断未遂。

\paragraph{边界与变式}
若甲不能游泳但能及时呼救，应分别验证各项措施，而不能认定所有 $C$ 都是假。若事故后的任何可行救助均已不能避免该死亡，$A$ 不成立，不能将死亡归给后续不救，先行致害行为仍独立评价。误以为落水者不是自己监护的幼儿，可能阻断对义务基础事实的故意认识。知道监护事实却误以为法律不要求施救，属于法律认识问题，应另查责任层面。若两项同时存在且同等紧迫的义务只能履行其一，须另查教材的义务冲突规则，不能从 $D$ 直接得出有责。

\section*{闭卷记忆与自测}

\paragraph{复述路线}
先写「形、控、险」，分别查外部形态、意识控制、新风险。有人参与风险时，加问「谁做、谁知、谁控」。见到不作为，把措施 $m$ 写在纸上，背「应、能、未、可、等」，逐项还原为 $D,C,N,A,E_j$。最后补结果、罪过和阻却事由。把「应」再展开为「管源、护人」，各举一个例子。

\paragraph{题一}
乙自行攀登，甲中途自愿接手绳索，乙已经悬空并依赖甲。甲在本可拉回时故意放弃继续救助。为何不能沿用最初的自陷风险结论？

\paragraph{核对}
至少要在接手前、接手后分别求 $S_y$ 和 $F_y$。接手后控制与依赖变更，可能产生 $D_{\text{人}}$。再看能否拉回、是否足以避免结果、最终实施的是积极松手还是仅不继续拉，以及相应罪过。初始自陷并不固定后续风险的归属。

\paragraph{题二}
肇事者有救助义务，现场却已无任何措施能避免伤者死亡。哪一项为假，能否推出先行驾驶行为无罪？

\paragraph{核对}
该死亡的 $A$ 为假，后续不作为的既遂归责链断开。先行驾驶另查违反注意义务、结果及归责，不能从后行为的结论倒推前行为无罪。
